\documentclass[11pt]{article}
\usepackage{tp2style}
\hypersetup{pdftitle={TP2 Team Norms},pdfauthor={CSE 360 Team 24}}
\fancyfoot[L]{\sffamily\scriptsize\color{muted}TP2 Team Norms \textperiodcentered{} Team 24}

\newcommand{\norm}[1]{\subsection*{#1}}

\begin{document}

\doctitle{TP2 Team Norms}{\team \\ Revised October 12, 2026, from the TP1 Team Norms of September 15, 2026}

\begin{minipage}{\linewidth}
\small
\textbf{How to read this document.} This is the TP1 text with every change marked.
\add{Blue underlined text} was added for TP2; \del{red struck-through text} was removed.
Each changed norm ends with a \textsf{\textbf{Why changed}} line saying what prompted the change.
Norms that read \textsf{No change} are word for word as agreed in TP1. Norms 11 and 12 are new.
\end{minipage}

\medskip
\begin{tabularx}{\linewidth}{@{}l X l@{}}
\toprule
\textbf{Norm} & \textbf{Change} & \textbf{Prompted by} \\
\midrule
1 & Meetings set by spacing, not fixed weekdays; summary presenter named in the meeting & TP1 standups \\
3 & Entity classes finished and tested before GUI work starts & TP1 integration \\
5 & Javadoc updated in the same commit as the code & TP1 review \\
6 & AI coaching added on top of validation & TP2 Task 5 \\
8 & ``Done'' requires a positive and a negative test and Javadoc & TP2 Tasks 3, 4, 7 \\
10 & Every member checks Canvas after each upload & TP2 instructions \\
11 & New: equal presenting time, two presenters, a named backup & TP1 recording \\
12 & New: pre-submission checklist & TP1 submission \\
\bottomrule
\end{tabularx}

\section*{Why we wrote these}
Each norm below deals with a specific problem that has caused team projects to fail before.
\add{After TP1 we reviewed each one and changed those that TP1 showed were not strong enough.}
We agreed to all of them together, and each of us agrees to follow them.

\section*{Our norms}

\norm{1. Meet for a standup at least twice a week}
We meet \del{Monday and Thursday, with more than 24 hours between meetings.}
\add{at least twice each week, with at least one full day between two meetings, at a time
agreed at the end of the previous standup. Before the meeting ends, one member is named to
produce its Status Summary screencast.} Each person takes two or three minutes to say what
they finished, what they will do next, and what is blocking them.

Why it matters: short, frequent meetings let us catch problems early. If two people are about
to change the same class, we find out now instead of during final integration.
\why{in TP1 we sometimes ended a standup without deciding who would make the summary video,
which delayed the videos. Fixed Monday and Thursday slots also cannot hold the extra standups
the TP2 schedule needs in its busiest week; the spacing rule keeps the gap the TP2 instructions
require between meetings.}

\norm{2. Record every standup to the Zoom cloud}
We record to the cloud and not to our own computers.

Why it matters: cloud recordings are compressed, so the summary videos stay well under the
100 MB limit. Local recordings are far too large.
\unchanged

\norm{3. Integrate our work continuously}
Nobody saves their work until the end. We merge small working pieces and test the whole
application together at least twice a week.
\add{A class that other code is built on is finished and passes its tests before that other
code is started. For TP2 this means LessonsLearned and LessonsLearnedList come before any
lesson window.}

Why it matters: splitting up the work, building alone, and combining everything at the end is
the most common way these projects fail.
\why{TP1 showed that a shared class is safest when it is finished and tested before the code
that depends on it is written. TP2 puts three new windows on two new entity classes.}

\norm{4. Use branches and review all code}
Every change goes on its own branch and needs one other person to read it before it is merged.

Why it matters: it stops two people from quietly changing the same method, and it catches
mistakes before they reach the shared code.
\unchanged

\norm{5. Keep the documents and the code matching}
When the code changes, the design documents and the comments change with it.
\add{A method's Javadoc is changed in the same commit as the method, and the reviewer checks
it. Screenshots of code are taken again whenever the code they show changes.}

Why it matters: if it is not written down, we have to keep explaining it to each other. The
documents and the code should look like one person wrote them.
\why{in TP1 our comments sometimes lagged behind the code they described. TP2 grades the
Javadoc and screenshots of the code directly, so a stale comment or screenshot now costs marks.}

\norm{6. Check every text input before using it}
Every text field has a maximum length. Fields that carry real meaning, such as username,
password, and email, get full checks.
\add{Lesson fields that benefit from it also get coaching (spelling, completeness, wording)
after they pass validation. Coaching advises and never blocks a save. The plan is in
AICoaching.pdf.}

Why it matters: typing a very long string into a field is a common way to crash a program. If
our program crashes, we have shown an attacker a weakness.
\why{TP2 Task 5 asks the team to plan for input coaching on top of the validation built in
TP1, and our HW2 study of LessonShare showed coaching at entry time improves what gets stored.}

\norm{7. Answer team messages within 12 hours, and report blockers right away}
If something is blocking us, we say so immediately instead of waiting for the next standup.

Why it matters: nobody can help with a problem they do not know about. Waiting a day wastes a
day.
\unchanged

\norm{8. Agree on what ``done'' means}
A story is done when it is written, tested, documented, reviewed, and merged.
\add{``Tested'' means at least one positive and one negative test that pass and are listed in
the Requirements to Tests Map. ``Documented'' means Javadoc on the class, every attribute, and
every method, including private ones.}

Why it matters: work that is ``almost done'' hides problems until it is too late to fix them.
\why{the TP1 definition did not say what ``tested'' or ``documented'' meant. TP2 Task 7
requires positive and negative tests for every requirement, and Tasks 3 and 4 grade the
documentation of each attribute.}

\norm{9. Deal with a teammate who is not helping}
We talk to the person directly and early. If nothing changes, we write down what happened and
tell the instructional team.

Why it matters: ignoring it rewards the person who is not working and lowers all of the
grades.
\unchanged

\norm{10. Upload a safety copy by noon on the due date}
We submit whatever we have at noon, then keep working and submit again.
\add{After every upload, each member opens the submission in Canvas, downloads it, and posts
in the team chat that it opens and is complete.}

Why it matters: the last successful submission is the one that is graded. A noon upload gives
us hours to fix any problem with the file or with Canvas.
\why{the TP2 instructions say every member must verify the submission in Canvas personally,
and that this cannot be delegated to one member.}

\norm{11. \add{Share presenting time equally}}
\add{Every main screencast has two presenters. Across the three main screencasts and the
summary screencasts, each member presents for roughly the same total time. A backup presenter
is named before each recording session.}
\why{(new norm) in TP1 one member was unavailable for a recording, and the member who stood in
had not prepared. TP2 Task 9 requires every member to present for roughly equal time.}

\norm{12. \add{Run a checklist before every upload}}
\add{Before the ZIP is made, one member who did not build the deliverable checks every file
and folder name against the assignment, compiles the project from a clean copy, runs every test
suite, deletes all .DS\_Store files, and confirms each video is MP4 or MOV and under 100 MB.}
\why{(new norm) the TP1 upload included 33 stale class files from an old build, and two standup
files used the wrong date format in their names. A file over 100 MB would make the whole TP2
submission score zero.}

\section*{How we handle disagreements}
Anyone can raise a problem privately first. If it keeps happening, we discuss it as a team at
the next standup. If it still is not fixed, we write it down and take it to the instructional
team.
\unchanged

\section*{Signatures}
By signing below, each of us agrees to follow these norms as revised for TP2.

\medskip
\begin{tabularx}{\linewidth}{@{}l X l@{}}
\toprule
\textbf{Member} & \textbf{Signature} & \textbf{Date} \\
\midrule
Arush Ugal & & \\[6pt]
Karan Patel & & \\[6pt]
Anirudh B & & \\[6pt]
Pranjal Ganvir & & \\[6pt]
Moksha Dev & & \\[6pt]
\bottomrule
\end{tabularx}

\end{document}
